\section{Der Kompilierungsprozess}
Il codice sorgente viene convertito in dati binari attraverso diverse fasi:

\begin{enumerate}
    \item \textbf{Preprocessore:} Elabora inclusioni e macro.
    (\texttt{\#define}, \texttt{\#if}, \dots)
    \item \textbf{Compilatore:} Converte il linguaggio di alto livello in assembly.
    \item \textbf{Assemblatore:} Converte l'assembly in codice macchina rilocabile.
    \item \textbf{Linker/Bindatore:} Collega i file oggetto e le librerie.
    \item \textbf{Loader:} Converte gli indirizzi rilocabili in indirizzi assoluti e li carica in memoria.
\end{enumerate}

\subsection{Compiler-Analyse}
\begin{itemize}
    \item \textbf{Analisi lessicale:} Rimozione degli spazi bianchi e raggruppamento dei token.
    \item \textbf{Analisi sintattica/semantica:} Verifica della struttura e dei tipi di dati.
    \item \textbf{Ottimizzazione:} Miglioramento del codice in termini di velocità o spazio di memoria.
    \item \textbf{Generazione:} Generazione del codice di destinazione.
\end{itemize}

\section{I/O-Schnittstellen}
I dispositivi periferici vengono mappati in memoria tramite registri specifici (\textit{Memory Mapped I/O}):
\begin{itemize}
    \item \textbf{Registri dati:} Contengono i dati da elaborare.
    \item \textbf{Registri di controllo:} Configurazione del dispositivo periferico.
    \item \textbf{Registri di stato:} Mostrano lo stato attuale dell'I/O (ad es. Busy, Ready).
\end{itemize}